\documentclass[11pt]{beamer}
\usetheme{Madrid}
\usecolortheme{default}

\setbeamercolor{structure}{fg=blue!70!black}
\setbeamercolor{block title}{bg=blue!70!black,fg=white}
\setbeamercolor{block body}{bg=blue!5,fg=black}
\setbeamercolor{block title alerted}{bg=red!70!black,fg=white}
\setbeamercolor{block body alerted}{bg=red!5,fg=black}
\setbeamertemplate{navigation symbols}{}
\setbeamertemplate{footline}[frame number]

\usepackage[utf8]{inputenc}
\usepackage[T1]{fontenc}
\usepackage[french]{babel}
\usepackage{tikz}
\usetikzlibrary{arrows.meta,positioning,shapes.geometric,fit,backgrounds}
\usepackage{booktabs}
\usepackage{array}
\usepackage{xcolor}

\title{De la maquette à une représentation hiérarchique}
\subtitle{Étude comparative de quatre approches de génération de code frontend}
\author{KOUONGA TADIE KENNEDY}
\institute{Université de Douala}
\date{}

\begin{document}

% ===================== TITLE =====================
\begin{frame}
\titlepage
\end{frame}

% ===================== PLAN =====================
\begin{frame}{Plan}
\begin{enumerate}
\item Question de recherche et périmètre de l'étude
\item Les quatre approches étudiées
\item Ce qui différencie leurs entrées
\item Comment mesurer la qualité d'une reconstruction
\item Comparaison chiffrée des quatre approches
\item Un mécanisme de correction automatique
\item Aspects pratiques : coût et bibliothèques de style
\item Limites actuelles et pistes de recherche
\end{enumerate}
\end{frame}

% =========================================================
\section{Question de recherche}

\begin{frame}{Question de recherche}
\begin{block}{Problème étudié}
Comment passer d'une représentation visuelle d'une interface (maquette, capture d'écran, croquis) à une \textbf{représentation structurée et hiérarchique} qui pourra ensuite servir à générer automatiquement du code frontend~?
\end{block}
\begin{center}
\begin{tikzpicture}[node distance=1.1cm, every node/.style={font=\small}]
\node[draw, rounded corners, fill=gray!10, minimum width=3.2cm, minimum height=1cm] (a) {Maquette / capture};
\node[draw, rounded corners, fill=blue!10, minimum width=3.2cm, minimum height=1cm, right=of a] (b) {Modèle hiérarchique};
\node[draw, rounded corners, fill=green!10, minimum width=3.2cm, minimum height=1cm, right=of b] (c) {Code frontend};
\draw[-{Latex}] (a) -- (b);
\draw[-{Latex}] (b) -- (c);
\end{tikzpicture}
\end{center}
\end{frame}

\begin{frame}{Périmètre exact de cette étude}
\begin{alertblock}{Où se situe le travail présenté}
Cette étude se concentre uniquement sur la \textbf{première flèche} du schéma précédent~: comment transformer la maquette en modèle hiérarchique (arbre de composants, arbre de mise en page).
\end{alertblock}
\begin{center}
\begin{tikzpicture}[node distance=1.1cm, every node/.style={font=\small}]
\node[draw, rounded corners, fill=gray!10, minimum width=3.2cm, minimum height=1cm] (a) {Maquette / capture};
\node[draw, rounded corners, fill=blue!20, thick, minimum width=3.2cm, minimum height=1cm, right=of a] (b) {Modèle hiérarchique};
\node[draw, rounded corners, fill=green!5, dashed, minimum width=3.2cm, minimum height=1cm, right=of b] (c) {Code frontend};
\draw[-{Latex}, thick] (a) -- (b);
\draw[-{Latex}, dashed, gray] (b) -- (c);
\node[below=0.2cm of a, font=\tiny] {étudié en détail};
\node[below=0.2cm of b, font=\tiny] {étudié en détail};
\node[below=0.2cm of c, font=\tiny, gray] {hors périmètre};
\end{tikzpicture}
\end{center}
La génération finale du code (deuxième flèche) n'est pas l'objet direct de la comparaison~: elle sert seulement à évaluer, \textit{a posteriori}, si le modèle hiérarchique construit était de bonne qualité.
\end{frame}

\begin{frame}{Pourquoi une représentation intermédiaire ?}
\begin{itemize}
\item Une image décrit surtout \textbf{ce que l'on voit}.
\item Le code doit exprimer \textbf{comment les éléments sont organisés}.
\item Il faut donc retrouver les éléments, leurs relations, leur organisation spatiale et leurs comportements de mise en page.
\end{itemize}
\begin{block}{Point important}
« Représentation hiérarchique » est utilisée au sens large~: arbre de composants, arbre de mise en page, arbre de subdivisions ou structure avec contraintes.
\end{block}
\end{frame}

% =========================================================
\section{Les quatre approches}

\begin{frame}{Les quatre approches étudiées}
\small
\begin{tabular}{p{3cm}p{8.5cm}}
\toprule
\textbf{Approche} & \textbf{Idée principale} \\
\midrule
Prototype2Code & Reconstruire une hiérarchie à partir des relations de contenance, puis organiser en lignes/colonnes \\
DCGen & Diviser récursivement l'image en segments ; la division forme un arbre \\
ScreenCoder & Détecter et nommer les grandes régions, puis construire un arbre selon des règles de conception web \\
DesignCoder & Diviser l'interface, extraire le sens des éléments, puis les regrouper dans un arbre de composants \\
\bottomrule
\end{tabular}
\end{frame}

\begin{frame}{Ce que les quatre méthodes ont en commun}
\begin{center}
\begin{tikzpicture}[node distance=0.9cm, every node/.style={font=\small}]
\node[draw, rounded corners, fill=blue!10, minimum width=2.6cm, minimum height=1cm] (d) {Détecter};
\node[draw, rounded corners, fill=blue!10, minimum width=2.6cm, minimum height=1cm, right=of d] (r) {Relier};
\node[draw, rounded corners, fill=blue!10, minimum width=2.6cm, minimum height=1cm, right=of r] (s) {Structurer};
\node[draw, rounded corners, fill=blue!10, minimum width=2.6cm, minimum height=1cm, right=of s] (g) {Générer};
\draw[-{Latex}] (d) -- (r);
\draw[-{Latex}] (r) -- (s);
\draw[-{Latex}] (s) -- (g);
\end{tikzpicture}
\end{center}
\begin{enumerate}
\item Identifier des unités de l'interface (éléments, régions, widgets)
\item Déterminer des relations entre ces unités
\item Construire une structure (ne plus traiter l'image comme plate)
\item Utiliser cette structure pour générer le code
\end{enumerate}
\end{frame}

% =========================================================
\section{Des entrées différentes}

\begin{frame}{Deux familles d'entrées}
Avant de comparer les résultats, il faut préciser ce que chaque méthode reçoit réellement en entrée~: cela change la difficulté du problème résolu.
\vspace{0.3cm}
\begin{columns}
\begin{column}{0.48\textwidth}
\begin{block}{Avec métadonnées de conception}
\small
\begin{itemize}
\item Prototype2Code
\item DesignCoder
\end{itemize}
\textit{Image + fichier Figma/Sketch (id, type, position, style déjà connus)}
\end{block}
\end{column}
\begin{column}{0.48\textwidth}
\begin{block}{Image seule}
\small
\begin{itemize}
\item DCGen
\item ScreenCoder
\end{itemize}
\textit{Aucune donnée cachée : tout doit être deviné depuis les pixels}
\end{block}
\end{column}
\end{columns}
\vspace{0.3cm}
\small\textit{Comme monter un meuble avec la notice de montage, contre le monter seulement à partir d'une photo du résultat fini.}
\end{frame}

% =========================================================
\section{Mesurer la qualité}

\begin{frame}{Comment mesurer si une reconstruction est bonne ?}
Deux questions distinctes se posent, et chacune a sa propre façon d'être mesurée~:
\begin{block}{1. Est-ce que ça \textit{ressemble} à l'original ?}
On compare les images (originale contre générée) --- similarité \textbf{visuelle}.
\end{block}
\begin{block}{2. Est-ce que c'est \textit{bien rangé} à l'intérieur ?}
On compare les arbres de composants (référence contre généré) --- similarité \textbf{structurelle}.
\end{block}
\small\textit{Deux sites peuvent se ressembler visuellement tout en ayant un code totalement désorganisé derrière --- d'où l'intérêt de mesurer les deux séparément.}
\end{frame}

\begin{frame}{Les métriques de similarité visuelle, expliquées simplement}
\begin{block}{CLIP score}
\small Un réseau de neurones compare l'image originale et l'image générée et donne un score entre 0 et 1~: plus le score est proche de 1, plus les deux images \og{}veulent dire la même chose\fg{} visuellement.
\end{block}
\begin{block}{Block-match / Texte / Position / Couleur}
\small On découpe les deux images en blocs (un bouton, un titre, une carte...), on les fait correspondre un à un, puis on vérifie si chaque bloc a la bonne taille, le bon texte, la bonne position et la bonne couleur.
\end{block}
\end{frame}

\begin{frame}{Les métriques de similarité structurelle, expliquées simplement}
\begin{block}{TreeBLEU}
\small Compare l'arbre du code généré à l'arbre de référence, morceau par morceau (un peu comme on compterait combien de groupes de mots se retrouvent à l'identique entre deux phrases).
\end{block}
\begin{block}{Container Match}
\small Vérifie si les éléments ont été regroupés dans les bons \og{}contenants\fg{} (le bon parent), sans juger les détails à l'intérieur --- juste~: est-ce que ce bouton a bien été rangé dans la bonne carte ?
\end{block}
\begin{block}{Tree Edit Distance (TED)}
\small Nombre minimum d'opérations (ajouter, supprimer, renommer un nœud) nécessaires pour transformer l'arbre généré en l'arbre de référence. \textbf{Plus ce nombre est bas, mieux c'est.}
\end{block}
\end{frame}

% =========================================================
\section{Comparaison chiffrée}

\begin{frame}{Fidélité visuelle : les scores CLIP}
\begin{center}
\small
\begin{tabular}{lc}
\toprule
\textbf{Méthode} & \textbf{Score CLIP (ressemblance visuelle)} \\
\midrule
ScreenCoder & 81 \% \\
DesignCoder & 92 \% \\
DCGen & 86 \% \\
Prototype2Code & non publié dans le document source \\
\bottomrule
\end{tabular}
\end{center}
\vspace{0.3cm}
\small\textit{Ces scores proviennent des évaluations propres à chaque article et permettent une première comparaison quantitative, à interpréter en tenant compte de la différence d'entrées vue plus haut.}
\end{frame}

\begin{frame}{Fidélité structurelle : au-delà de l'apparence}
\begin{center}
\small
\begin{tabular}{lccc}
\toprule
\textbf{Méthode} & \textbf{TreeBLEU} & \textbf{Container Match} & \textbf{TED} \\
\midrule
Prototype2Code & 0.53 & 0.58 & 40.0 \\
DesignCoder & \textbf{0.69} & \textbf{0.75} & \textbf{28.2} \\
\bottomrule
\end{tabular}
\end{center}
\vspace{0.3cm}
\small\textit{DesignCoder obtient un score supérieur sur les trois métriques structurelles~: son arbre de composants est non seulement visuellement fidèle, mais aussi mieux organisé --- notamment parce qu'il cherche à comprendre le \textbf{rôle} de chaque élément avant de le ranger.}
\end{frame}

% =========================================================
\section{Correction automatique}

\begin{frame}{Un mécanisme de correction automatique}
Une des quatre méthodes ne se contente pas de produire le code en une seule passe~: elle vérifie et corrige son propre résultat.
\begin{center}
\begin{tikzpicture}[node distance=0.6cm, every node/.style={font=\footnotesize}]
\node[draw, rounded corners, align=center, fill=blue!10, minimum width=2.3cm, minimum height=0.9cm] (g) {Code généré};
\node[draw, rounded corners, align=center, fill=blue!10, minimum width=2.3cm, minimum height=0.9cm, right=of g] (r) {Rendu \\ photographié};
\node[draw, rounded corners, align=center, fill=blue!10, minimum width=2.3cm, minimum height=0.9cm, right=of r] (c) {Comparaison \\ à l'original};
\node[draw, rounded corners, align=center, fill=green!10, minimum width=2.3cm, minimum height=0.9cm, right=of c] (f) {Correction \\ automatique};
\draw[-{Latex}] (g) -- (r);
\draw[-{Latex}] (r) -- (c);
\draw[-{Latex}] (c) -- (f);
\draw[-{Latex}] (f) to[bend left=40] (g);
\end{tikzpicture}
\end{center}
\small\textit{Sans ce module, les scores de DesignCoder chutent nettement (CLIP~: 0.92 $\to$ 0.88 ; TreeBLEU~: 0.66 $\to$ 0.55) --- ce qui confirme que cette boucle de vérification a un effet mesurable sur la qualité finale.}
\end{frame}

% =========================================================
\section{Aspects pratiques}

\begin{frame}{Coût et bibliothèques de style}
\begin{block}{Bibliothèques de composants}
ScreenCoder génère directement en \textbf{Tailwind CSS}, une bibliothèque de styles standard, plutôt que du CSS entièrement fait main.
\end{block}
\begin{block}{Temps de génération}
ScreenCoder discute explicitement le compromis \textbf{coût contre qualité}~: un modèle plus petit permet un rendu rapide et approximatif, un modèle plus grand prend plus de temps (de l'ordre de quelques dizaines de secondes) pour un résultat plus fidèle.
\end{block}
\begin{alertblock}{Ce qui reste ouvert}
Aucune des quatre méthodes ne gère de design system personnalisé (tokens de couleur ou de typographie propres à une marque).
\end{alertblock}
\end{frame}

% =========================================================
\section{Limites et pistes}

\begin{frame}{Limite : une image figée dans le temps}
Les quatre méthodes construisent un arbre à partir d'une image \textbf{statique}, à un instant donné.
\begin{itemize}
\item Aucune ne modélise le survol de la souris, l'ouverture d'un menu ou d'une fenêtre modale.
\item Aucune ne modélise l'adaptation dynamique à différentes tailles d'écran, au-delà de la mise en page initiale.
\end{itemize}
\end{frame}

\begin{frame}{Limite : l'accessibilité n'est pas prise en compte}
Aucune des quatre méthodes ne mentionne les normes d'accessibilité~:
\begin{itemize}
\item Pas de rôles sémantiques HTML5 profonds (\texttt{main}, \texttt{nav}, \texttt{article}).
\item Pas de prise en compte des lecteurs d'écran ni de la navigation au clavier.
\end{itemize}
\small\textit{Ce constat vaut pour l'ensemble des quatre approches, pas seulement pour certaines d'entre elles.}
\end{frame}

\begin{frame}{Limite : dépendance à un langage cible unique}
\begin{center}
\small
\begin{tabular}{ll}
\toprule
\textbf{Méthode} & \textbf{Technologie cible unique} \\
\midrule
DesignCoder & React Native uniquement \\
Prototype2Code, DCGen, ScreenCoder & HTML/CSS uniquement \\
\bottomrule
\end{tabular}
\end{center}
\vspace{0.3cm}
Aucune méthode ne produit un format intermédiaire réellement indépendant de la technologie de sortie.
\end{frame}

\begin{frame}{Conclusion et pistes de recherche}
\begin{block}{Synthèse}
Les travaux étudiés convergent vers une même idée~: la maquette doit être transformée en une représentation structurée avant de produire le code. La qualité de cette représentation intermédiaire conditionne directement la qualité du code final.
\end{block}
\begin{alertblock}{Question qui se dégage}
Quel niveau d'information doit contenir notre modèle pour être indépendant de la technologie cible tout en restant suffisamment riche pour générer ensuite React, HTML/CSS, Flutter, etc. ?
\end{alertblock}
\begin{block}{Trois pistes concrètes}
\begin{enumerate}
\item Un format intermédiaire réellement indépendant du langage cible
\item Une gestion des états d'interface (survol, clic, responsive)
\item Une prise en compte de l'accessibilité dès la construction de l'arbre
\end{enumerate}
\end{block}
\end{frame}

\begin{frame}{Références principales}
\footnotesize
\begin{itemize}
\item S. Xiao et al. \textit{Prototype2Code}. arXiv:2405.04975, 2024.
\item Y. Wan et al. \textit{Divide-and-Conquer: Generating UI Code from Screenshots}. Proc. ACM Softw. Eng., FSE, 2025.
\item Y. Jiang et al. \textit{ScreenCoder}. arXiv:2507.22827, 2025.
\item Y. Chen et al. \textit{DesignCoder: Hierarchy-Aware and Self-Correcting UI Code Generation}. arXiv:2506.13663, 2025.
\end{itemize}
\end{frame}

\end{document}
